\corrigesChapitre{11}{Temps d’attente et lois à densité}
\begin{corrige}{1}{Premier succès}
\begin{questions}
\item $0{,}8^2\times0{,}2=0{,}128$.
\item $0{,}8^4=0{,}4096$.
\item $E(X)=1/0{,}2=5$ essais.
\end{questions}
\end{corrige}
\begin{corrige}{2}{Une uniforme}
\begin{questions}
\item $1/6$ sur $[2,8]$, et zéro ailleurs.
\item $(5-3)/(8-2)=1/3$.
\item $E(X)=5$, $V(X)=6^2/12=3$.
\end{questions}
\end{corrige}
\begin{corrige}{3}{Une exponentielle}
\begin{questions}
\item $\mathrm e^{-1}\simeq0{,}367879$.
\item $1-\mathrm e^{-1}\simeq0{,}632121$.
\item $1/0{,}5=2$ minutes.
\end{questions}
\end{corrige}
\begin{corrige}{4}{Hauteur et aire}
\begin{questions}
\item $1/0{,}2=5$.
\item C’est l’aire totale qui doit valoir $1$ : ici $5\times0{,}2=1$. Une hauteur n’est pas une probabilité.
\item $0$, car une loi à densité attribue une probabilité nulle à un point isolé.
\end{questions}
\end{corrige}
\begin{corrige}{5}{Vérifier une densité}
\begin{questions}
\item $\int_0^2kx\,\mathrm dx=2k=1$ donne $k=1/2$ ; la fonction est bien positive.
\item $\int_0^1x/2\,\mathrm dx=1/4$.
\item $\int_0^2x^2/2\,\mathrm dx=4/3$.
\end{questions}
\end{corrige}
\begin{corrige}{6}{Fonction de répartition}
\begin{questions}
\item $F(x)=0$ si $x<0$, $x^2$ si $0\leq x\leq1$, et $1$ si $x>1$.
\item $(3/4)^2-(1/2)^2=5/16$.
\item $E(X)=2/3$, $E(X^2)=1/2$, donc $V(X)=1/2-4/9=1/18$.
\end{questions}
\end{corrige}
\begin{corrige}{7}{Attendre encore}
\begin{questions}
\item $0{,}7^7/0{,}7^4=0{,}7^3=0{,}343$.
\item Les deux probabilités sont égales.
\item Après quatre échecs, la probabilité d’attendre encore plus de trois essais est identique à celle du départ ; les essais futurs conservent le même modèle indépendant.
\end{questions}
\end{corrige}
\begin{corrige}{8}{Retrouver le paramètre}
\begin{questions}
\item $\mathrm e^{-10\lambda}=1/2$, donc $\lambda=\ln2/10$.
\item $10/\ln2\simeq14{,}427$ unités de temps.
\item Par absence de mémoire, $P(T>10)=1/2$.
\end{questions}
\end{corrige}
\begin{corrige}{9}{Bus et mémoire}
\begin{questions}
\item $6/12=1/2$.
\item $P(U>9)/P(U>3)=(3/12)/(9/12)=1/3$.
\item Non : l’absence de mémoire imposerait ici $P(U>6)=1/2$, différent de $1/3$.
\end{questions}
\end{corrige}
\begin{corrige}{10}{Simuler une exponentielle}
\begin{questions}
\item $T\leq t$ équivaut à $\ln(1-U)\geq-\lambda t$, donc $U\leq1-\mathrm e^{-\lambda t}$.
\item L’uniformité donne $P(T\leq t)=1-\mathrm e^{-\lambda t}$ : c’est l’exponentielle.
\item $1-U$ est strictement positif si $U\in[0,1[$ ; le logarithme est alors défini, même lorsque $U=0$.
\end{questions}
\end{corrige}
\begin{corrige}{11}{Discrétiser une attente}
\begin{questions}
\item $P(K>n)=P(T>nh)=\mathrm e^{-\lambda nh}$.
\item C’est une géométrique sur $\N^*$ avec $1-p=\mathrm e^{-\lambda h}$, donc $p=1-\mathrm e^{-\lambda h}$.
\item L’approximation locale $\mathrm e^{-z}\simeq1-z$ donne $p\simeq\lambda h$ pour $\lambda h$ petit ; ce n’est pas une égalité générale.
\end{questions}
\end{corrige}
\begin{corrige}{12}{Une densité impossible}
\begin{questions}
\item $\int_0^1(3x-1)\,\mathrm dx=1/2$.
\item $f<0$ pour $0\leq x<1/3$.
\item Non : une multiplication positive conserve les valeurs négatives ; une multiplication négative rend négative l’autre partie. L’intégrale égale à $1$ ne suffirait pas.
\end{questions}
\end{corrige}
